\par\noindent\textbf{(1)} 定义 \(L_n=\int d\varepsilon\,\Sigma(\varepsilon)(\varepsilon-\mu)^n[-\partial_\varepsilon f]\)，\(\Sigma(\varepsilon)\) 为能量分辨电导核。在线性响应中 \(\sigma=e^2L_0\)，电子热导 \(\kappa_e=(L_2-L_1^2/L_0)/T\)，Seebeck系数 \(S=-L_1/(eTL_0)\)。
\par\noindent\textbf{(2)} 在 \(\Sigma\) 平滑条件下，\(L_0=\Sigma(E_F)+O(T^2)\)，\(L_1=(\pi^2/3)(k_BT)^2\Sigma'(E_F)+O(T^4)\)，\(L_2=(\pi^2/3)(k_BT)^2\Sigma(E_F)+O(T^4)\)。奇偶性保证常数项对 \(L_1\) 不贡献。
\par\noindent\textbf{(3)} 最低阶有 \(\kappa_e/T\approx(\pi^2/3)k_B^2\Sigma(E_F)\)，与 \(\sigma=e^2\Sigma(E_F)\) 相除得 \(L_0^{\mathrm{Lorenz}}=\kappa_e/(\sigma T)=\pi^2k_B^2/(3e^2)\)。数值约 \(2.44\times10^{-8}\,\mathrm{W\,\mathrm{ohm}\,K^{-2}}\)，在弹性散射主导简并金属中成立。
\par\noindent\textbf{(4)} 修正项 \(L_1^2/L_0\sim T^4\) 对 \(\kappa_e\) 为 \(O(T^3)\)，因此不影响领先Lorenz数。强非弹性散射、非费米液体、极低载流子浓度或声子热导占优时，\(\kappa/(\sigma T)\) 可显著偏离常数；不能将这个结果当作所有材料的精确恒等式。
